\documentclass[10pt]{article}
\usepackage[T1]{fontenc}
\usepackage[utf8]{inputenc}
\usepackage[french]{babel}
\usepackage{amsmath,amssymb,mathrsfs,mathpazo}
\usepackage{geometry}
\geometry{a4paper,margin=2cm}
\usepackage{graphicx,array,tabularx,tabularray,colortbl,multirow,multicol,booktabs,enumitem,pifont,ulem,fancybox,diagbox,cancel}
\usepackage{pstricks,pstricks-add,pst-plot,pst-tree,pst-node,pst-text,pst-func,pst-math,pst-bspline,pst-3dplot,pst-eucl,pst-solides3d}
\let\pscalculate\undefined
\usepackage{pst-bezier}
\usepackage{tikz,pgfplots,tkz-tab}
\usetikzlibrary{arrows,arrows.meta,patterns,calc,positioning,trees,decorations.markings,shapes.geometric,shapes.misc,shapes.arrows,intersections,decorations.pathmorphing,decorations.pathreplacing,backgrounds,shadows,babel}
\pgfplotsset{compat=1.18}
\usepgfplotslibrary{fillbetween}
\usepackage[np]{numprint}
\usepackage{listings,makecell,pas-tableur,twoopt}
\let\barre\undefined
\usepackage[pstricks]{tablvar}
\definecolor{darkgreen}{RGB}{0,100,0}
\providecommand{\red}{\color{red}}
\providecommand{\blue}{\color{blue}}
\providecommand{\green}{\color{green}}
\providecommand{\black}{\color{black}}
\providecommand{\white}{\color{white}}
\providecommand{\gray}{\color{gray}}
\providecommand{\cyan}{\color{cyan}}
\providecommand{\magenta}{\color{magenta}}
\providecommand{\yellow}{\color{yellow}}
\providecommand{\violet}{\color{violet}}
\providecommand{\orange}{\color{orange}}
\providecommand{\euro}{\texteuro}
\providecommand{\EUR}{\texteuro}
\providecommand{\ieme}{\textsuperscript{e}}
\providecommand{\degres}{\ensuremath{^\circ}}
\providecommand{\degree}{\ensuremath{^\circ}}
\providecommand{\vv}[1]{\overrightarrow{#1}}
\providecommand{\Coord}[1]{#1}
\providecommand{\decosix}{}
\providecommand{\starredbullet}{\ensuremath{\bullet}}
\pagestyle{empty}
\setlength{\parindent}{0pt}
\usepackage{adjustbox}\DeclareUnicodeCharacter{2236}{:}\DeclareUnicodeCharacter{2190}{\ensuremath{\leftarrow}}
\providecommand{\R}{\mathbb{R}}
\providecommand{\N}{\mathbb{N}}
\providecommand{\Z}{\mathbb{Z}}
\providecommand{\Q}{\mathbb{Q}}
\providecommand{\D}{\mathbb{D}}
\providecommand{\C}{\mathbb{C}}
\providecommand{\e}{{\mathrm{e}}}
\providecommand{\vect}[1]{\overrightarrow{#1}}
\providecommand{\vectt}[1]{\overrightarrow{\mathrm{#1}}}
\providecommand{\vvt}[1]{\overrightarrow{\mathrm{#1}}}
\providecommand{\np}[2][]{{#2~#1}}
\providecommand{\vv}[1]{{\overrightarrow{#1}}}
\renewcommand{\d}{{\mathrm{d}}}
\providecommand{\degres}{{^\circ}}
\providecommand{\degre}{{^\circ}}
\providecommand{\blue}{\color{blue}}
\providecommand{\red}{\color{red}}
\providecommand{\enskip}{\quad}
\providecommand{\llbracket}{\lbrack\!\lbrack}
\providecommand{\rrbracket}{\rbrack\!\rbrack}
\providecommand{\euro}{{\text{€}}}
\providecommand{\no}{{\text{n°}}}
\providecommand{\interval}[2]{{[#1;#2]}}
\providecommand{\intervFO}[2]{{[#1;#2[}}
\providecommand{\dotfill}{\ldots}

\providecommand{\ds}{\displaystyle}
\providecommand{\textding}[1]{#1}
\providecommand{\Oij}{\(\mathrm{O};\vec{\imath},\vec{\jmath}\)}
\providecommand{\Oijk}{\(\mathrm{O};\vec{\imath},\vec{\jmath},\vec{k}\)}
\providecommand{\Ouv}{\(\mathrm{O};\vec{u},\vec{v}\)}
\providecommand{\intervFF}[2]{\ensuremath{[#1;#2]}}
\renewcommand{\euro}{\ifmmode\text{\texteuro}\else\texteuro\fi}
\DeclareUnicodeCharacter{00B0}{\ifmmode^{\circ}\else\textdegree\fi}
\DeclareUnicodeCharacter{20AC}{\ifmmode\text{\texteuro}\else\texteuro\fi}
% Énoncés : Éducation nationale. Transcription et corrigés : APMEP, Denis Vergès et contributeurs. Découpage et classement : CoopMaths, Sébastien Lozano, Éric Elter et contributeurs. Adaptation de la présentation : Maths Explorer.
% CC BY-SA : https://coopmaths.fr/www/about/
% https://coopmaths.fr/alea/static/bac/2022/tex/bac_2022_05_sujet1_asie_2.tex
\begin{document}
\noindent\textit{APMEP / CoopMaths. CC BY-SA. Présentation adaptée par Maths Explorer.}\par\medskip
\medskip

\emph{Principaux domaines abordés} : Suites numériques. Algorithmique et programmation.

\medskip

Un médicament est administré à un patient par voie intraveineuse. 

\bigskip

\textbf{Partie A : modèle discret de la quantité médicamenteuse}

\medskip

Après une première injection de 1 mg de médicament, le patient est placé sous perfusion.

On estime que, toutes les $30$~minutes, l'organisme du patient élimine 10\,\% de la quantité de médicament présente dans le sang et qu'il reçoit une dose supplémentaire de $0,25$ mg de la substance médicamenteuse.

On étudie l'évolution de la quantité de médicament dans le sang avec le modèle suivant :

pour tout entier naturel $n$, on note $u_n$ la quantité, en mg, de médicament dans le sang du patient au bout de $n$ périodes de trente minutes. On a donc $u_0 = 1$.

\medskip

\begin{enumerate}
\item Calculer la quantité de médicament dans le sang au bout d'une demi-heure.
\item Justifier que, pour tout entier naturel $n,\: u_{n+1} = 0,9u_n + 0,25$.
\item 
	\begin{enumerate}
		\item Montrer par récurrence sur $n$ que, pour tout entier naturel $n,\: u_n \leqslant  u_{n+1} < 5$.
		\item En déduire que la suite $\left(u_n\right)$ est convergente.
	\end{enumerate}
\item On estime que le médicament est réellement efficace lorsque sa quantité dans le sang
du patient est supérieure ou égale à $1,8$ mg.
	\begin{enumerate}
		\item Recopier et compléter le script écrit en langage Python suivant de manière à déterminer au bout de combien de périodes de trente minutes le médicament commence à être réellement efficace.
\begin{center}
\begin{tabular}{|l|}\hline
\texttt{\textbf{def} efficace():}\\
\quad \texttt{u=1}\\
\quad \texttt{n=0}\\
\quad \texttt{\textbf{while} \ldots\ldots :}\\
\quad\qquad \texttt{u=\ldots\ldots}\\
\quad\qquad \texttt{n = n+1}\\
\quad \texttt{\textbf{ return }n}\\ \hline
\end{tabular}
\end{center}

		\item Quelle est la valeur renvoyée par ce script ? Interpréter ce résultat dans le contexte de l'exercice.
	\end{enumerate}		
\item Soit $\left(v_n\right)$ la suite définie, pour tout entier naturel $n$, par $v_n = 2,5 - u_n$.
	\begin{enumerate}
		\item Montrer que $\left(v_n\right)$ est une suite géométrique dont on précisera la raison et le premier terme $\left(v_0\right)$.
		\item Montrer que, pour tout entier naturel $n,\: u_n = 2,5 - 1,5 \times 0,9^n$.
		\item Le médicament devient toxique lorsque sa quantité présente dans le sang du patient dépasse $3$ mg. 
		
D'après le modèle choisi, le traitement présente-t-il un risque pour le patient? Justifier.
	\end{enumerate}
\end{enumerate}

\bigskip

\textbf{Partie B : modèle continu de la quantité médicamenteuse}

\medskip

Après une injection initiale de $1$~mg de médicament, le patient est placé sous perfusion.

Le débit de la substance médicamenteuse administrée au patient est de $0,5$ mg par heure.

La quantité de médicament dans le sang du patient, en fonction du temps, est modélisée
par la fonction $f$, définie sur $[0~;~ +\infty[$, par 

\[f(t) = 2,5 - 1,5\text{e}^{-0,2t},\]

où $t$ désigne la durée de la perfusion exprimée en heure.

On rappelle que ce médicament est réellement efficace lorsque sa quantité dans le sang du patient est supérieure ou égale à $1,8$~mg.

\medskip

\begin{enumerate}
\item Le médicament est-il réellement efficace au bout de 3~h 45~min ?
\item Selon ce modèle, déterminer au bout de combien de temps le médicament devient réellement efficace.
\item Comparer le résultat obtenu avec celui obtenu à la question 4. b. du modèle discret de la Partie A.
\end{enumerate}

\bigskip

\end{document}
